\documentclass[../main]{subfiles}
\begin{document}

\chapter{Teorema de Bernoulli}

\section{Teorema de Bernoulli}

\subsection{Principios y aplicaciones del Teorema de Bernoulli}

El Teorema de Bernoulli es una de las leyes fundamentales de la dinámica de fluidos. Fue formulado por el matemático y físico suizo Daniel Bernoulli en el siglo XVIII. Este teorema describe el comportamiento de un fluido en movimiento a lo largo de una corriente de flujo, y tiene numerosas aplicaciones en la ingeniería y la física, especialmente en la mecánica de fluidos.

\info{Teorema de Bernoulli}{
El Teorema de Bernoulli establece que en una corriente de flujo estacionario, la suma de la presión estática, la presión dinámica y la presión debida a la altura es constante a lo largo de una línea de corriente. Matemáticamente, se expresa como:
\begin{equation}
P + \frac{1}{2} \rho v^2 + \rho g h = \text{constante}
\end{equation}

Siendo:
\begin{itemize}
    \item $P$: Presión estática del fluido [Pa]
    \item $\rho$: Densidad del fluido [kg/m³]
    \item $v$: Velocidad del fluido [m/s]
    \item $g$: Aceleración de la gravedad [m/s²]
    \item $h$: Altura respecto a un nivel de referencia [m]
\end{itemize}

Esta ecuación se aplica a fluidos incomprensibles y no viscosos en flujo laminar.
}

Una de las aplicaciones más comunes del Teorema de Bernoulli es en el diseño de aviones. Las alas de los aviones están diseñadas de tal manera que la velocidad del aire es mayor en la parte superior que en la inferior. Esto genera una menor presión en la parte superior del ala, creando una fuerza de sustentación que mantiene el avión en el aire.

Otra aplicación importante es en el diseño de sistemas de tuberías. En un sistema de tuberías, el teorema puede utilizarse para calcular las velocidades y presiones del fluido en diferentes puntos de la tubería, asegurando que el sistema funcione de manera eficiente y segura.

\actividad{Responder el siguiente cuestionario}
\begin{enumerate}
    \item ¿Qué establece el Teorema de Bernoulli?
    \item ¿Cuáles son las tres presiones que se suman en la ecuación de Bernoulli?
    \item ¿Cómo afecta la velocidad del fluido a la presión estática según el Teorema de Bernoulli?
    \item ¿Qué suposiciones se hacen sobre el fluido en el Teorema de Bernoulli?
    \item ¿Qué aplicaciones prácticas tiene el Teorema de Bernoulli en la ingeniería?
    \item ¿Cómo se relaciona el Teorema de Bernoulli con la sustentación en las alas de un avión?
    \item ¿Qué es la presión dinámica y cómo se calcula?
    \item ¿Por qué el Teorema de Bernoulli no se aplica a fluidos viscosos?
    \item ¿Cómo puede el Teorema de Bernoulli ayudar a diseñar sistemas de tuberías eficientes?
    \item ¿Cuál es la importancia de la densidad del fluido en la ecuación de Bernoulli?
\end{enumerate}

\actividad{Resolver los siguientes problemas}
\begin{enumerate}
    \item Calcular la presión estática $P$ en un punto de una tubería donde el agua (densidad $\rho = 1000$ kg/m³) fluye a una velocidad $v = 5$ m/s y la altura $h$ es 10 m. La presión en otro punto a la misma altura, donde la velocidad es 3 m/s, es $50000$ Pa.
    \item Encontrar la velocidad del fluido $v$ en una sección de la tubería si en otra sección la velocidad es 2 m/s, la presión es 30000 Pa y la presión en la primera sección es 20000 Pa. La densidad del fluido es $1000$ kg/m³ y ambas secciones están al mismo nivel.
    \item Determinar la altura $h$ a la que se eleva el fluido si la presión en la base es $75000$ Pa, la velocidad del fluido es 0 m/s en la base, y la densidad del fluido es $1000$ kg/m³. La presión en el punto más alto es $50000$ Pa.
    \item Calcular la densidad $\rho$ de un fluido si la presión en un punto es $60000$ Pa, la velocidad en ese punto es 4 m/s, y a otro punto a la misma altura, donde la velocidad es 2 m/s, la presión es $80000$ Pa.
    \item Encontrar la presión $P$ en una sección estrecha de un tubo horizontal si en una sección más ancha la presión es $100000$ Pa, la velocidad en la sección estrecha es 8 m/s y en la sección más ancha es 2 m/s. La densidad del fluido es $1000$ kg/m³.
    \item Determinar la velocidad $v$ en la parte más baja de una tubería vertical si la presión en la parte superior es $40000$ Pa, la velocidad en la parte superior es 1 m/s, y la altura entre ambos puntos es 5 m. La densidad del fluido es $1000$ kg/m³.
    \item Calcular la diferencia de presión $\Delta P$ entre dos puntos de una tubería horizontal si la velocidad en el primer punto es 3 m/s y en el segundo punto es 6 m/s. La densidad del fluido es $1000$ kg/m³.
    \item Encontrar la altura $h$ en un tubo vertical si en la base la velocidad del fluido es 0 m/s, la presión es $50000$ Pa y en el punto más alto la presión es $25000$ Pa. La densidad del fluido es $1000$ kg/m³.
    \item Calcular la presión $P$ en un punto de una tubería donde la velocidad del fluido es 3 m/s, la densidad del fluido es $1000$ kg/m³, y la presión en otro punto a la misma altura y con velocidad 1 m/s es $40000$ Pa.
    \item Determinar la velocidad $v$ del fluido en una sección de una tubería horizontal si en otra sección la velocidad es 4 m/s y la diferencia de presión entre las dos secciones es $15000$ Pa. La densidad del fluido es $1000$ kg/m³.
\end{enumerate}

\end{document}
